% !TeX root = ../../Book3.tex
% 纲要标题：### 18.4 正则局部环的阶乘性
% 纲要位置：计划纲要.md，第 2070 行。

\section{正则局部环的阶乘性}
\label{b3:sec:1804}

% TODO: 按以下纲要要点撰写本节。

% 纲要内容（仅作写作计划，尚非教材正文）：
%
% * 证明正则局部环是唯一分解整环，先把目标化为高度一素理想的主性；完整保留 Noether 局部假设
% * 准备所需行列式线引理：利用有限自由消解控制主开集上的可逆模；在使用处明确有限展示模延拓与局部化正合性的论证
% * 选择 \(x\in\mathfrak m\setminus\mathfrak m^2\)，结合 \(R/(x)\) 的正则性、维数归纳及主开集上的秩一模分析，完成高度一素理想主性的证明
% * 区分“每个素点附近局部为主”与“给定主开集上的可逆模为自由”两个步骤；不能只凭局部阶乘性推断全局唯一分解
% * 将第一卷唯一分解、第二卷行列式线与本卷局部同调结构连成完整证明链；第20章据此讨论除子与局部阶乘性。证明方法核对见 [V13-4]
%

\subsection{唯一分解与高度一素理想的主性}
\label{b3:subsec:1804:01}

待撰写。

\subsection{行列式线与主开集上的可逆模}
\label{b3:subsec:1804:02}

待撰写。

\subsection{正则参数、维数归纳与秩一模}
\label{b3:subsec:1804:03}

待撰写。

\subsection{局部主性与主开集上的自由性}
\label{b3:subsec:1804:04}

待撰写。

\subsection{正则局部环阶乘性的证明及其与除子的联系}
\label{b3:subsec:1804:05}

待撰写。
